% !PRINTED 624
\weberpage{624}

\[
\begin{array}{lcccccc}
 & a_1\sqrt{\Delta}, & a_2, & a_3\sqrt{\Delta}, & a_4, & a_5\sqrt{\Delta}, & a_6\\
\text{die Grade} & 2 & 4 & 6 & 8 & 10 & 12.
\end{array}
\]
$\sqrt{\Delta}$ ist aber vom 10$^{\text{ten}}$ Grade in den $x$, und folglich muss $a_1=0$,
$a_3=0$ und $a_5$ eine Zahl sein.

Man kann die Co\"efficienten $a_2$, $a_4$, $a_5$, $a_6$ durch wirkliche
Bildung der Ausdr\"ucke nach den Vorschriften \"uber die Darstellung der symmetrischen
Functionen berechnen, was von \gesperrt{Cayley} geschehen ist. Rechnung und Formeln sind aber sehr
weitl\"aufig.

Wir wollen uns hier damit begn\"ugen, die Resolvente f\"ur
einen besonderen Fall zu bilden. Der dabei gefundene Werth
f\"ur die Zahl $a_5$ ist dann nat\"urlich allgemein g\"ultig.

Die Gleichung 5$^{\text{ten}}$ Grades habe die \gesperrt{Bring-Jerrard'sche}
Form\footnote{\gesperrt{Runge}, Acta mathematica, Bd.\,7.}
\weq{13}{x^5+\alpha x+\beta=0.}
Dann sind $a_2$, $a_4$, $a_6$ ganze rationale Functionen von $\alpha$ und $\beta$.
Von den Co\"efficienten der allgemeinen Gleichung (6) ist $d$ von
der 4$^{\text{ten}}$, $e$ von der 5$^{\text{ten}}$ Ordnung in den $x$. Daraus folgt, dass
in $a_2$, $a_4$, $a_6$ der Co\"efficient $e$ nur mit einem der Co\"efficienten
$a$, $b$, $c$ multiplicirt vorkommen kann, und dass also, wenn $a$, $b$, $c$
Null gesetzt werden, $e$ aus diesen Ausdr\"ucken herausgehen muss.
F\"ur die Gleichung (13) k\"onnen also $a_2$, $a_4$, $a_6$ nicht von $\beta$ ab-
h\"angen, und man erh\"alt mit R\"ucksicht auf den Grad, wenn
$m$, $m_1$, $m_2$, $m_3$ Zahlen bedeuten
\weq{14}{a_2=m_1\alpha,\quad a_4=m_2\alpha^2,\quad a_6=m_3\alpha^3,\quad a_5=m.}

Die Zahlen $m$, $m_1$, $m_2$, $m_3$ lassen sich aus einer speciellen
Annahme bestimmen. Setzen wir $\beta=0$, so wird
\weq{15}{\begin{aligned}
x_0&=0,\quad x_1=\sqrt[4]{-\alpha},\quad x_2=i\sqrt[4]{-\alpha},\\
&\hphantom{{}=0,\quad{}} x_3=-\sqrt[4]{-\alpha},\quad x_4=-i\sqrt[4]{-\alpha},
\end{aligned}}
und daraus ergiebt sich
\weq{16}{\begin{aligned}
\sqrt{\Delta} &= (x_0-x_1)(x_0-x_2)(x_0-x_3)(x_0-x_4)(x_1-x_2)\\
&\quad (x_1-x_3)(x_1-x_4)(x_2-x_3)(x_2-x_4)(x_3-x_4)\\
&= 16i\sqrt{-\alpha}^5 = -16\sqrt{\alpha}^5\\
\Delta &= 256\alpha^5.
\end{aligned}}
